% Reproducible spectral numerics and declared resolution controls.
% This is a fragment included by the manuscript; it intentionally has no preamble.

\section{Thickness-resolved spectral numerics and error control}
\label{sec:spectral-numerics}

This section fixes the discrete weak form, branch tracking, perturbation
solves, critical-point resolution checks, response quadrature, and acceptance
rules used by every numerical artifact.  The objective is reproducibility:
each symbol is tied below to a code interface, a qualified configuration
field, or a qualified artifact field.

\subsection{GLL element map and traction-free weak form}

We use the same nondimensionalization as in
Sec.~\ref{sec:isotropic-rayleigh-lamb-derivation}: \(h\) is the half-thickness,
\(z\in[-h,h]\), \(k=h k^{\rm phys}\), and
\(\omega=h\omega^{\rm phys}/c_T\).  A separated displacement has the form

\[
 \boldsymbol u(x,y,z,t)=\boldsymbol U(z)
 \exp\{\mathrm i(k_x x+k_y y)-\mathrm i\omega t\}.
\]

On element \(e=[z_e^-,z_e^+]\), let \(p\) be the polynomial degree and
\(\{\xi_j,w_j\}_{j=0}^{p}\) the \(p+1\) Legendre--Gauss--Lobatto (GLL)
nodes and weights returned by \texttt{gll\_nodes\_weights}.  The physical
map, Jacobian, weights, and derivative matrix are

\begin{equation}
 \begin{aligned}
 z_e(\xi)&=\frac{z_e^-+z_e^+}{2}
            +\frac{z_e^+-z_e^-}{2}\xi,\qquad -1\leq\xi\leq1,\\
 J_e&=\frac{z_e^+-z_e^-}{2},\qquad
 w_{ej}=J_e w_j,\qquad D_e=J_e^{-1}D .
 \end{aligned}
 \label{eq:gll-element-map}
\end{equation}

Here \(D\) is the barycentric matrix produced by
\texttt{differentiation\_matrix}.

The routine \texttt{build\_gll\_mesh}
shares the interface node of adjacent elements and maps the local
connectivities to node-major triples
\((U_x,U_y,U_z)\).  Thus one element on \([-h,h]\) is the production
discretization for a homogeneous plate, while a split mesh with bounds
\((-h,0,h)\) supplies an independent assembly check.

For test function \(\boldsymbol V\), integration by parts in \(z\) gives

\begin{equation}
 \int_{-h}^{h}\!\boldsymbol\varepsilon(\boldsymbol V)^*
 :\mathsf C:\boldsymbol\varepsilon(\boldsymbol U)\,dz
 =\omega^2\int_{-h}^{h}\!\rho\,\boldsymbol V^*\!\cdot
 \boldsymbol U\,dz
 +\sum_{z\in\{-h,h\}}
 \boldsymbol V(z)^*\!\cdot[\boldsymbol\sigma(z)\boldsymbol n(z)] .
 \label{eq:spectral-weak-form}
\end{equation}

No displacement degree of freedom is constrained at either face.  The
boundary term vanishes precisely when

\begin{equation}
 \bigl.\boldsymbol\sigma\boldsymbol n\bigr|_{z=\pm h}=0
 \quad\text{is the natural traction-free boundary condition.}
 \label{eq:natural-traction-boundary}
\end{equation}

\begin{equation}
 \boldsymbol\varepsilon_{\mathrm M}
 =(\varepsilon_{xx},\varepsilon_{yy},\varepsilon_{zz},
 \sqrt2\varepsilon_{yz},\sqrt2\varepsilon_{xz},
 \sqrt2\varepsilon_{xy})^{\mathsf T}.
 \label{eq:mandel-convention}
\end{equation}

This Mandel convention makes the Euclidean product equal to the tensor
contraction, \(\boldsymbol\varepsilon_{\mathrm M}^*
\mathsf C_{\mathrm M}\boldsymbol\varepsilon_{\mathrm M}\).  With nodal
interpolant \(\mathsf N_j\), the strain matrix at a quadrature node is
\(\mathsf B=\mathsf B_0+k_x\mathsf B_x+k_y\mathsf B_y\).  The element
and assembled matrices used by \texttt{assemble\_plate\_matrices} are

\[
 \begin{aligned}
 \mathsf K_e&=\sum_{j=0}^{p}w_{ej}\mathsf B_j^*
 \mathsf C_{\mathrm M}\mathsf B_j,\\
 \mathsf M_e&=\sum_{j=0}^{p}w_{ej}\rho\mathsf N_j^*\mathsf N_j,\\
 \mathsf K&=\mathop{\mathrm{assemble}}_e\mathsf K_e,\qquad
 \mathsf M=\mathop{\mathrm{assemble}}_e\mathsf M_e .
 \end{aligned}
\]

Consistent collocated GLL quadrature is used for stiffness and mass.  Because
the cardinal basis is evaluated at the same GLL nodes, the assembled mass is
diagonal.  This is quadrature-induced mass lumping; no separate
post-assembly lumping or penalty term is applied.  The reported Hermitian defect is
\(\|\mathsf K-\mathsf K^*\|_F/\|\mathsf K\|_F\).  Constitutive input is
rejected unless its minor and major symmetries hold and its Mandel matrix is
positive definite.

\subsection{Matrix derivatives and eigensystem diagnostics}

Because \(\mathsf B\) is affine in \((k_x,k_y)\), the five derivative
matrices are assembled analytically at the same quadrature nodes by the
wavevector-derivative assembler:

\begin{equation}
 \begin{aligned}
 \mathsf K_i&=\sum_{e,j}w_{ej}
 (\mathsf B_i^*\mathsf C_{\mathrm M}\mathsf B
 +\mathsf B^*\mathsf C_{\mathrm M}\mathsf B_i),\\
 \mathsf K_{ii}&=2\sum_{e,j}w_{ej}
 \mathsf B_i^*\mathsf C_{\mathrm M}\mathsf B_i,\\
 \mathsf K_{xy}&=\sum_{e,j}w_{ej}
 (\mathsf B_x^*\mathsf C_{\mathrm M}\mathsf B_y
 +\mathsf B_y^*\mathsf C_{\mathrm M}\mathsf B_x),
 \qquad i\in\{x,y\},\quad \partial_{k_i}\mathsf M=0.
 \end{aligned}
 \label{eq:wavevector-matrix-derivatives}
\end{equation}

The generalized Hermitian problem
\(\mathsf K\boldsymbol u_n=\lambda_n\mathsf M\boldsymbol u_n\),
\(\lambda_n=\omega_n^2\), is solved by \texttt{solve\_plate\_modes}.
The solver uses SciPy's generalized Hermitian \texttt{gvx} driver for the
lowest twelve modes and requires a positive-definite mass.  It rejects
\(\lambda_{\min}<-\tau_-\), while values in \([-\tau_-,0]\) are clipped to
zero, where
\[
 \tau_-=10^3\epsilon_{\rm mach}
 \max\!\left(\frac{\|\mathsf K\|_2}{\lambda_{\min}(\mathsf M)},
 \operatorname{tiny}_{64}\right).
\]
Every eigenvector is explicitly mass normalized.
\(\operatorname{tiny}_{64}\) denotes the smallest positive normal binary64
number.  For the selected mode the code's normalized residual is

\begin{equation}
 r=\frac{\|\mathsf Ku-\lambda\mathsf Mu\|_2}
 {\max\{\|\mathsf Ku\|_2+\|\lambda\mathsf Mu\|_2,
 \operatorname{tiny}_{64}\}}
 \quad\text{(normalized residual)}.
 \label{eq:normalized-eigen-residual}
\end{equation}

\begin{equation}
 d_M=\|\mathsf U^*\mathsf M\mathsf U-\mathsf I\|_F
 \quad\text{(mass orthogonality defect)}.
 \label{eq:mass-orthogonality-defect}
\end{equation}

\begin{equation}
 \operatorname{MAC}_M(u,v)
 =\frac{|u^*\mathsf Mv|^2}{(u^*\mathsf Mu)(v^*\mathsf Mv)}.
 \label{eq:mass-mac}
\end{equation}

\paragraph{Mode tracking.}
For two nonzero vectors, \texttt{mass\_mac} evaluates
Eq.~\eqref{eq:mass-mac}.  The continuation rule is
\begin{equation}
 \boxed{\begin{aligned}
 &\text{predict the eigenvalue and maximize the previous-subspace projection;}\\
 &\text{form the connected close-eigenvalue cluster and test its overlap;}\\
 &\text{require scalar mass MAC or squared cluster overlap, then phase align;}\\
 &\text{retain both the selected vector and the complete cluster basis.}
 \end{aligned}}
 \label{alg:mode-tracking}
\end{equation}

If \(s_{\min}\) is the smallest principal cosine between the previous and
current mass-orthonormal cluster bases, a step with threshold \(\eta\) passes
when
\[
 \operatorname{MAC}_M\geq\eta
 \quad\text{or}\quad s_{\min}^{\,2}\geq\eta .
\]
The scalar/squared-subspace threshold is \(\eta=0.2\) in the generic tracker
and \(\eta=0.8\) in the registered robustness sweeps.  The connected-cluster
relative eigenvalue threshold is \(10^{-3}\).

\texttt{RingAnchoredSpectralEvaluator} constructs eight
endpoint-free angular anchors independently at orders \(p\) and \(p+4\),
requires angular closure scalar MAC and the unsquared principal cosine each
to be at least \(1-10^{-8}\),
then continues monotonically away from \(q=0\) on separate inward and outward
rays.  It also rejects a cluster touching the top of the twelve computed
modes.  The relative eigengap returned by the tracker is

\begin{equation}
 \begin{aligned}
 g_{\rm rel}
 &=\min_{m\ne n}\frac{|\lambda_m-\lambda_n|}
 {\max(|\lambda_n|,|\lambda_m|,1)},\\
 g_{\rm rel}&>10\,\delta\lambda_{\rm rel}.
 \end{aligned}
 \label{eq:eigengap-rejection}
\end{equation}

A sample that fails the second line is rejected by the eigengap gate.

Here
\[
 \delta\lambda_{\rm rel}=
 \max\!\left\{
 \frac{|\lambda^{(p+4)}-\lambda^{(p)}|}
 {\max(|\lambda^{(p+4)}|,|\lambda^{(p)}|,1)},
 r^{(p)},r^{(p+4)}
 \right\}.
\]
The same nested-order pair defines absolute frequency and gradient discrepancy
estimators as the differences between the two tracked answers.  These
inter-order estimators quantify resolution sensitivity; without a continuum
enclosure or a posteriori theorem, they are not upper bounds on the unknown
continuum errors.
This gate is essential: every derivative formula below assumes a simple,
spectrally isolated branch.

\subsection{Sensitivities and critical-point resolution checks}

Let \(r\) denote the radial wavenumber and let \(\epsilon\) multiply the
cubic constitutive perturbation at fixed density.  At unit mass norm,
\(\lambda_p=u^*\mathsf K_pu\).  The radial mode derivative is found by the
bordered reduced-resolvent system implemented in
\texttt{radial\_frequency\_sensitivity}:

\begin{equation}
 (\mathsf K-\lambda\mathsf M)u_p
 =-(\mathsf K_p-\lambda_p\mathsf M)u,
 \qquad u^*\mathsf Mu_p=0,\qquad
 \lambda_p=u^*\mathsf K_pu .
 \label{eq:differentiated-mode-solve}
\end{equation}

Row and constraint scalings are applied before the dense solve; the original
equation and gauge are then checked with relative residual \(10^{-8}\).  The
mixed sensitivity includes both the direct matrix derivative and the mode
derivative,

\[
 \lambda_{r\epsilon}
 =u^*\mathsf K_{r\epsilon}u
 +2\operatorname{Re}(u^*\mathsf K_\epsilon u_r),\qquad
 V=\frac{\lambda_\epsilon}{2\omega},\qquad
 B=\frac{\lambda_{r\epsilon}}{2\omega}
 -\frac{\lambda_\epsilon\lambda_r}{4\omega^3}.
\]

For the primary sensitivity artifact, the matrix perturbations are constructed
at each angle from the exactly linear cubic family as
\[
 \mathsf K_\epsilon=
 \frac{\mathsf K(+0.25)-\mathsf K(-0.25)}{0.5},\qquad
 \mathsf K_{r\epsilon}=
 \frac{\mathsf K_r(+0.25)-\mathsf K_r(-0.25)}{0.5}.
\]
The full profile uses order \(10\), 64 endpoint-free angles, and a single
centered physical check with \(\Delta\epsilon=\Delta q=0.0025\); the smoke
profile uses order \(8\), 32 angles, and \(0.005\).  On an \(N_\theta\)-point
grid the reported cubic coefficient is the discrete Fourier coefficient
\[
 V_4=\frac{2}{N_\theta}\sum_{j=0}^{N_\theta-1}
 V(\theta_j)\cos(4\theta_j),
\]
and \(V_0=N_\theta^{-1}\sum_jV(\theta_j)\).  The primary artifact accepts
only when the relative vector errors of both \(V\) and \(B\), compared with
that single centered full-wave check, are below the configured field
\begin{center}
\texttt{config/reference.yaml::sensitivity\_match\_tolerance}.
\end{center}

The separate convergence artifact supplies step-refinement evidence.  It uses
the tensor-product centered grid
\(\Delta\epsilon,\Delta q\in\{0.02,0.01,0.005,0.0025\}\), checks the last two
rows/columns against the same tolerance, and requires improvement relative to
the coarsest row/column.  Thus only this independent sweep varies both step
sizes; it does not differentiate the perturbation implementation against
itself.

\begin{equation}
 \boxed{\begin{aligned}
 &\text{seed every sign-changing polar cell;}\\
 &\text{solve }\nabla_{\boldsymbol k}\omega=0\text{ and deduplicate;}\\
 &\text{classify the Cartesian Hessian;}\\
 &\text{check both annular boundaries and the index sum.}
 \end{aligned}}
 \label{alg:critical-point-search}
\end{equation}

The production implementation is the registered critical-point locator.  Its
search annulus is
\[
 A=\{\boldsymbol k:(1-f_A)k_0\leq|\boldsymbol k|
 \leq(1+f_A)k_0\}.
\]
Here \(f_A\) is \texttt{config/reference.yaml::annulus\_fraction}.  On a polar
candidate grid, the Cartesian gradient is resolved into radial and tangential
components.  Every cell across which both components bracket zero supplies a
seed.  SciPy's \texttt{scipy.optimize.root} follows each seed with its default
Powell hybrid method (\texttt{method=hybr}).  The finite-difference Cartesian
Hessian is supplied as the Jacobian, with solver tolerance
\(10^{-11}\).  Candidates outside \(A\) or with gradient residual larger
than their nested-order gradient-discrepancy estimator are rejected.  An unresolved
Hessian is not silently discarded: classification raises an error and fails
the resolution-check stage.  Remaining roots are merged only when their
separation does not exceed \(d_{\rm merge}\), where
\[
 \begin{aligned}
 b_0&=5\times10^{-9}\max(k_0,1),\\
 b_{\max}&=\max\{b_0,\min(0.05\,\Delta_A,0.01\,k_0)\},\\
 b_{12}&=\max\{b_0,\delta k_{{\rm lin},1}+\delta k_{{\rm lin},2}\},\\
 d_{\rm merge}&=\min(b_{12},b_{\max}),
 \end{aligned}
\]
where \(\Delta_A\) is the annular half-width.

The full workflow uses a \(5\times16\) candidate grid and independently
repeats the search on \(9\times32\).  Counts, kinds, positions, frequencies,
and Hessian eigenvalues must agree within the following propagated bounds.
For matched coarse/fine roots, define
\[
 \begin{aligned}
 d_k&=\|\boldsymbol k_c-\boldsymbol k_f\|_2,\\
 u_k&=\frac{\delta g_c}{\min_\ell|\lambda_\ell(H_c)|}
     +\frac{\delta g_f}{\min_\ell|\lambda_\ell(H_f)|},\\
 H_{\max}&=\max_{\nu\in\{c,f\}}\max_\ell|\lambda_\ell(H_\nu)|,\\
 d_H&=\max_\ell|\lambda_\ell(H_c)-\lambda_\ell(H_f)|.
 \end{aligned}
\]
The comparison passes only if
\[
 \begin{aligned}
 d_k&\leq\max(5\times10^{-9}k_0,u_k),\\
 |\omega_c-\omega_f|
 &\leq\delta\omega_c+\delta\omega_f
       +\tfrac12 H_{\max}d_k^2,\\
 d_H&\leq\delta H_c+\delta H_f.
 \end{aligned}
\]
These candidate-grid bounds are distinct from the capped deduplication radius
used within either search.  The
full critical-point profile uses Hessian step \(10^{-3}\); the smoke profile
uses \(2\times10^{-3}\).  Both use \(N_b=32\) and \(2N_b=64\) boundary
samples for the final annular consistency check.  The Green workflow repeats
the accepted search on the \(15\times32\) response-support domain specified
below; the silicon workflow uses \(9\times32\).  Both use boundary levels
\(16/32\); their Hessian steps are \(10^{-3}\) (full Green) and
\(3\times10^{-3}\) (silicon), respectively.  The
negative-\(\epsilon\) search is additionally rotated by \(\pi/16\), so its
success is not caused by placing candidate cells on the positive-\(\epsilon\)
symmetry rays.  The recorded position scale follows from the linearized
gradient equation:

\begin{equation}
 \delta k_{\rm lin}:=
 \frac{\delta g}{\min_\ell|\lambda_\ell(H)|}
 =\|H^{-1}\|_2\,\delta g,
 \qquad H=\nabla_{\boldsymbol k}^2\omega .
 \label{eq:critical-point-uncertainty}
\end{equation}

Equation~\eqref{eq:critical-point-uncertainty} is a local resolution scale,
not a nonlinear interval enclosure.  Its interpretation requires a
nondegenerate Hessian that remains stable over the displacement.  Here
\(\delta g\) is the code field \texttt{gradient\_uncertainty}.  The root
residual is checked separately and must not exceed \(\delta g\); a
conservative first-order displacement estimate under the linearized Hessian
would replace \(\delta g\) by
\(\delta g+\|\nabla_{\boldsymbol k}\omega\|_2\).

The Hessian is a Richardson extrapolation of centered differences of the
full-wave gradient.  Its uncertainty is the fine-versus-Richardson
difference plus the two nested-order gradient uncertainties propagated
through the stencil.  Classification is accepted only if the smallest
absolute Hessian eigenvalue exceeds ten times this uncertainty.  A minimum
or maximum has gradient index \(+1\), whereas a saddle has gradient index
\(-1\).  Despite the legacy artifact field name
\texttt{critical\_points.morse\_index}, it stores this Poincar\'e gradient
index, not the Morse index (the number of negative Hessian eigenvalues).
The pointwise \texttt{gradient\_uncertainty} drives acceptance and uncertainty
propagation but is not stored as a critical-point array; the sidecar stores
the boundary maxima and the artifact stores
\texttt{critical\_points.gradient\_residual}.

\begin{equation}
 \operatorname{wind}_{\partial A}(\nabla_k\omega)
 =\sum_{k_j\in A}\operatorname{ind}(k_j),
 \quad\text{and the boundary winding index closes.}
 \label{eq:annular-index-certificate}
\end{equation}

The independent routine \texttt{verify\_annular\_exhaustion} samples both
annular boundaries with \(N_b\) and \(2N_b\) endpoint-free angles.  It
requires (i) the minimum boundary-gradient norm to exceed ten times the
maximum gradient uncertainty, (ii) equal coarse and fine boundary winding,
(iii) every fine angular increment of the gradient direction to be less than
\(\pi/2\), and (iv) the outer-minus-inner boundary winding to equal the sum
of accepted Poincar\'e indices.  Thus the boundary winding and the root index
closes on both resolutions.  This is a finite-resolution boundary/index
consistency check: by itself it cannot exclude an unresolved, mutually canceling
critical-point pair in the annular interior.  Candidate-grid doubling,
symmetry-offset repetition, and Hessian separation provide the additional
acceptance checks used here.  The perturbative eight-point prediction is
stored for comparison, but its prediction error is not a root-acceptance
gate.

\subsection{Polar Green quadrature, fit masks, and phase accuracy}

The endpoint-free grid is \(\theta_j=2\pi j/N_\theta\).  The production Green
workflow exploits verified \(C_{4v}\) covariance: its
\texttt{\_CachedCubicSymmetryFactory} maps every angle to
\[
 r_\theta=\theta\bmod\frac{\pi}{2},\qquad
 \theta_{\rm can}=\min\!\left(r_\theta,\frac{\pi}{2}-r_\theta\right)
\]
and rounds \(\theta_{\rm can}\) and the radial cache key to 14 decimal
places.  It reuses samples at equal
\((\theta_{\rm can},\operatorname{sign}q,k_0+q)\).  Each canonical ray starts
fresh inward and outward ring-anchored spectral continuations at \(q=0\).
Symmetry-equivalent rays therefore reuse the
canonical full-wave value; they are not recomputed as independent solves.
Every refinement level owns a separate cache, so no coarse-grid result is
used as a fine-grid answer.

At a surface node the phase-invariant normal-impulse modal factor is

\[
 A(q,\theta)=\frac{\mathrm i}{2\omega(q,\theta)}
 \left|u_z(z=h;q,\theta)\right|^2 .
\]

Let \(A^{(p)}\) and \(A^{(p+4)}\) denote the two spectral-order amplitudes.
The full-wave radial-evaluator factory requires
\[
 |A^{(p+4)}-A^{(p)}|
 \leq10^{-12}+10^{-2}\max(|A^{(p+4)}|,|A^{(p)}|).
\]
At \(q=0\), independently created inward and outward anchors must also agree:
\[
 \begin{aligned}
 |\omega_{\rm in}-\omega_{\rm out}|
 &\leq\delta\omega_{\rm in}+\delta\omega_{\rm out}
 +32\epsilon_{\rm mach}\max(\omega_{\rm in},\omega_{\rm out},1),\\
 |A_{\rm in}-A_{\rm out}|
 &\leq\delta A_{\rm in}+\delta A_{\rm out}
 +32\epsilon_{\rm mach}\max(|A_{\rm in}|,|A_{\rm out}|,1).
 \end{aligned}
\]

The Gaussian source and the local super-Gaussian branch window are fixed,
not fitted:

\[
 S(k)=\exp[-(Rk/2)^2],\qquad
 W(q)=\exp[-(q/\sigma)^8]b_\sigma(q).
\]

The ratio \(R/h\) is mapped to
\begin{center}
\texttt{config/reference.yaml::source\_radius\_over\_h}.
\end{center}
The ratio \(\sigma/k_0\) is mapped to
\begin{center}
\texttt{config/reference.yaml::window\_sigma\_over\_k0}.
\end{center}

The deterministic \(C^\infty\) taper \(b_\sigma\) equals one for
\(|q|\leq1.25\sigma\), follows the standard flat \(\exp(-1/x)\) transition,
and vanishes identically for \(|q|\geq1.5\sigma\).  The radial interval
\([-1.5\sigma,1.5\sigma]\) therefore contains the complete support, contains
exactly one \(q=0\) node, keeps \(k_0+q>0\), and has no endpoint term.  The
endpoint-free trapezoidal rule in
\(\theta\) and composite Simpson rule in \(q\) evaluate the Cartesian
Fourier measure without using a Bessel approximation:

\begin{equation}
 \begin{aligned}
 F_{\ell j}&=(k_0+q_\ell)S(k_0+q_\ell)W(q_\ell)A_{\ell j},\\
 G_h(t)&=\frac{\operatorname{Simpson}_{q_\ell}}{(2\pi)^2}
 \left[\frac{2\pi}{N_\theta}
 \sum_{j=0}^{N_\theta-1}F_{\ell j}
 e^{-\mathrm i\omega_{\ell j}t}\right],\\
 \widetilde G_h(t)&=e^{\mathrm i\omega_{\rm c}t}G_h(t),
 \qquad \omega_{\rm c}=\omega_0+\epsilon V_0 .
 \end{aligned}
 \label{eq:polar-green-quadrature}
\end{equation}

The registered-grid integration interface exposes the generic three-level
construction.  The production Green workflow rebuilds
each level as an independent full-wave tracked surface, while retaining exact
nested nodes for discrepancy evaluation.  It uses two distinct registered
three-grid sequences:

\[
 \begin{array}{c|ccc}
 \text{uniform-crossover rows }(\epsilon\ne0.08)
   &129\times32&257\times64&513\times128\\
 \text{fixed-anisotropy row }(\epsilon=0.08)
   &129\times64&257\times128&513\times256
 \end{array}
\]

The fixed-Morse comparison searches the tracked branch on
\(|q|\leq1.55\sigma=0.2325k_0\), which strictly contains the direct support
\(|q|\leq1.5\sigma=0.225k_0\) with margin \(0.05\sigma\).  It uses a
\(15\times32\) candidate grid and 16/32-node boundary/index checks.  Every
resolved stationary point in this search domain is multiplied by the same
\(b_\sigma\) and included in the Morse sum.  The sidecar records the two
domains, margin, boundary checks, point counts, and contribution counts for
each \(\epsilon\).

The dimensions are radial-by-angular.  Separate caches are created at every
level, so refinement does not reuse a coarse numerical frequency as a fine
answer.  By contrast, the full-profile convergence artifact's auxiliary
robustness study constructs a \(513\times128\) reference surface and uses
nested subsampling to \(65\times16\), \(129\times32\), and
\(257\times64\).  The smoke profile uses a \(257\times128\) reference and
\(33\times16\), \(65\times32\), and \(129\times64\) subsamples.  This
auxiliary study measures quadrature and interpolation convergence; it is not
substituted for the independent full-wave production grids.

The auxiliary responses use the same 61 times
\(t_m=1500+75m\), \(0\leq m\leq60\).  If \(G_\ell\) is the response from
subsampled level \(\ell\) and \(G_{\rm ref}\) is the reference response, then
\[
 \|f\|_{t,{\rm RMS}}=
 \left[M_t^{-1}\sum_m|f(t_m)|^2\right]^{1/2},
 \qquad
 Q_\ell=\frac{\|G_\ell-G_{\rm ref}\|_{t,{\rm RMS}}}
 {\|G_{\rm ref}\|_{t,{\rm RMS}}}.
\]
The stored \(Q_\ell\) field is
\texttt{convergence.quadrature\_error}.
To measure frequency interpolation, each level first receives a periodic
cubic spline in \(\theta\), evaluated at every reference angle; an ordinary
cubic spline in \(q\) is then evaluated at every reference radial node.
Writing the resulting surface as
\(\mathcal I_q\mathcal I_\theta\omega_\ell\), the stored error is
\[
 \begin{aligned}
 I_\ell&=\max_{(q,\theta)\in\mathcal G_{\rm ref}}
 |e_\ell(q,\theta)|,\\
 e_\ell&=\mathcal I_q\mathcal I_\theta\omega_\ell-\omega_{\rm ref}.
 \end{aligned}
\]
The stored field is \texttt{convergence.interpolation\_error}.
Thus the angular-periodic interpolation precedes the radial interpolation;
the order is part of the method.

Let \(G_0,G_1,G_2\) be the demodulated responses on three successive
production grids.  Their RMS complex differences
\(E_j=\|G_{j+1}-G_j\|_{\rm RMS}\) must decrease unless the fine discrepancy
is below the roundoff floor
\[
 E_{\rm round}=64\epsilon_{\rm mach}
 \max(\|G_2\|_{\rm RMS},1).
\]
The final discrepancy must satisfy

\[
 E_1\leq10^{-8}+0.05\|G_2\|_{\rm RMS},\qquad
 E_1/E_0\leq0.5
\]

whenever the rate lies above the absolute/roundoff floor.  The auxiliary
surface study separately requires decreasing quadrature and interpolation
errors, refinement ratios at most \(0.5\), and a finest relative quadrature
error below \(0.05\).

The decay exponents are extracted only from preregistered masks.  With
\(\tau=\epsilon V_4t\),

\begin{equation}
 \begin{aligned}
 M_{\rm early}(\epsilon,t)
 &=\mathbf 1_{\{\epsilon=\epsilon_{\rm first}\}}
   \mathbf 1_{\{t\geq1500\}}
   \mathbf 1_{\{0.10\leq|\tau|\leq0.30\}},\\
 M_{\rm late}(\epsilon,t)
 &=\mathbf 1_{\{\epsilon=0.08\}}
   \sum_{b}\mathbf 1_{[t_b,t_b+T_{\rm ms})}(t),\\
 T_{\rm ms}&=2\pi/|\omega_{\rm saddle}-\omega_{\rm min}|,
 \qquad t_0=1500 .
 \end{aligned}
 \label{eq:fit-masks}
\end{equation}

These two arrays are stored as
\texttt{green\_crossover.fit\_window\_early} and
\texttt{green\_crossover.fit\_window\_late}.
Here \(\epsilon_{\rm first}=0.005\) in the committed full profile and
\(\epsilon_{\rm first}=0.01\) in the smoke profile; operationally the code
uses \texttt{row\_index == 0}.  The early mask must contain at least ten
positive-amplitude samples.  The
call \texttt{fit\_power\_law} performs ordinary least squares on
\((\log t,\log|G|)\) and fits both the diagnostic decay slope and a nuisance
log--log intercept.  That intercept is used only for slope estimation.  The
separate full-wave/asymptotic comparison applies no fitted alignment of
amplitude, phase, carrier frequency, or time.  For the
late mask, each beat bin is half-open, contains at least
three time samples, and is retained only when the complete interval lies in
the computed time axis.  The fit uses at least four such bins, their geometric
time centers, and the direct-response RMS in each bin.  The independent
fixed-Morse comparison uses \(1500\leq t\leq10200\), excluding only samples
whose exact-Morse phasor coherence is below
\texttt{\_MORSE\_COHERENCE\_THRESHOLD} \(=0.30\).  Uniform collapse is checked
only for \(\epsilon\leq0.04\), \(t\geq1500\), and
\(|\tau|\leq\texttt{\_UNIFORM\_TAU\_MAXIMUM}=2.0\); the
\(\epsilon=0.08\) fixed-Morse row is excluded from that first-order gate.
The fixed-Morse RMS and maximum discrepancies are normalized by the
exact-Morse RMS over the retained noncancellation mask; the cancellation-region
RMS uses the same denominator.  The uniform-collapse diagnostic is the
maximum complex absolute error between
\texttt{green\_crossover.scaled\_response} and
\texttt{green\_crossover.J0} on its registered mask.

The operational crossover is the first downward crossing of scaled-response
magnitude \(0.9\), linearly interpolated between adjacent time samples within
\([0.5t_{\rm pred},1.5t_{\rm pred}]\).  Absence of such a crossing fails the
stage.  Because the search bracket is centred on the theoretical
\(t_{\rm pred}\propto|\epsilon V_4|^{-1}\), the fitted log slope of crossover
time versus \(\epsilon\) is a theory-centred consistency diagnostic rather
than independent evidence for inverse-rate scaling.  The complex Bessel
collapse is the primary numerical test.

Two production accumulated inter-resolution phase-discrepancy estimators are
enforced as acceptance tests.  The direct full-wave estimator combines (i)
each surface's nested spectral order-\(p\)/order-\(p+4\) frequency difference
and (ii) the common-node frequency difference between adjacent polar grids.
The exact-Morse estimator uses the largest nested-order frequency difference
among its eight resolved critical-point contributions.  Both call
\texttt{assert\_phase\_accuracy}, which enforces

\begin{equation}
 t_{\max}\max|\delta\omega|\le0.05 .
 \label{eq:phase-error-certificate}
\end{equation}

The right-hand side is the configured field
\begin{center}
\texttt{config/reference.yaml::phase\_error\_tolerance}.
\end{center}
For a grid pair the
\(\max|\delta\omega|\) in Eq.~\eqref{eq:phase-error-certificate} is the
maximum of its common-node discrepancy and the two adjacent surfaces' largest
local spectral discrepancies.  The stored
\texttt{green\_crossover.phase\_error} is the larger of the direct-grid and
exact-Morse accumulated discrepancy estimators.  Its legacy field name does
not turn this quantity into a continuum error bound.

The auxiliary convergence artifact uses a distinct interpolation discrepancy
estimator,
\[
 \Delta\phi_\ell=t_{\max}
 \left(\delta\omega_{{\rm interp},\ell}
 +\max\delta\omega_{{\rm local},\ell}
 +\max\delta\omega_{{\rm local},{\rm ref}}\right).
\]
Only its finest level is gated against the same \(0.05\) budget.  This
diagnostic is not substituted for either production estimator.
In short, the declared comparison has no fitted alignment and no
response-derived fitting window.

\subsection{Declared configuration and fixed acceptance thresholds}

The following list is the complete live top-level configuration registry.
No numerical value used by the workflows is silently supplied by this
document.

{\small
\begin{itemize}
\item \texttt{config/reference.yaml::schema\_version}: contract version one.
\item \texttt{config/reference.yaml::h}: half-thickness \(h=1\).
\item \texttt{config/reference.yaml::rho}: density \(\rho=1\).
\item \texttt{config/reference.yaml::lambda}:\\
Lam\'e modulus \(\lambda=2\), loaded as \texttt{ReferenceConfig.lam}.
\item \texttt{config/reference.yaml::mu}: shear modulus \(\mu=1\).
\item \texttt{config/reference.yaml::delta}: unit cubic perturbation amplitude.
\item \texttt{config/reference.yaml::epsilon\_values}: ordered nonzero signed
perturbations from \(-0.08\) to \(0.08\).
\item \texttt{config/reference.yaml::source\_radius\_over\_h}: \(R/h=0.5\).
\item \texttt{config/reference.yaml::window\_sigma\_over\_k0}:
\(\sigma/k_0=0.15\).
\item \texttt{config/reference.yaml::window\_sensitivity}: widths \(0.10\)
and \(0.20\) for robustness reruns.
\item \texttt{config/reference.yaml::annulus\_fraction}: \(f_A=0.15\).
\item \texttt{config/reference.yaml::zgv\_search\_kappa}:
\(0.20\leq\kappa\leq1.40\).
\item \texttt{config/reference.yaml::zgv\_search\_omega}:
\(2.40\leq\omega\leq3.60\).
\item \texttt{config/reference.yaml::eigen\_residual\_tolerance}: \(10^{-10}\).
\item \texttt{config/reference.yaml::isotropic\_match\_tolerance}: \(10^{-7}\).
\item \texttt{config/reference.yaml::curvature\_match\_tolerance}: \(10^{-4}\).
\item \texttt{config/reference.yaml::sensitivity\_match\_tolerance}: \(10^{-4}\).
\item \texttt{config/reference.yaml::phase\_error\_tolerance}: \(0.05\) rad.
\end{itemize}
}

Additional fixed pass/fail thresholds are part of the versioned
implementation rather than the YAML file.  The complete production list is:
{\small
\begin{enumerate}
\item The generalized matrices must be Hermitian to relative \(10^{-12}\);
the convergence sweep requires Hermitian defect below \(10^{-12}\), mass
orthogonality defect below \(10^{-10}\), and normalized eigen residual below
the configured eigen-residual tolerance.  Final relative
\(k_0,\omega_0\) errors must each be below the isotropic-match tolerance and
the curvature error below the curvature-match tolerance.

\item Every selected eigengap must exceed ten times its stated numerical
uncertainty.  An angular anchor's order-\(p\)/order-\(p+4\) eigenvalue
discrepancy must also be less than \(0.1\) times the resolved gap, angular
closure MAC and principal cosine must each be at least \(1-10^{-8}\), and a
selected cluster may not touch the upper end of the twelve-mode solve.
The independent branch sweep additionally requires reference-to-tracked MAC
at least \(0.99\) and gap larger than ten times its maximum eigen residual.

\item The differentiated-mode equation and gauge residual must be at most
\(10^{-8}\).  Primary relative \(V\) and \(B\) finite-difference errors, and
the last-two-step \(V_4\) and \(B\) convergence errors, must be below the
sensitivity-match tolerance.  The step sweep must improve from its coarsest
row and column.

\item A scaling-stage radial root must have residual no larger than
\[
 \max\!\left\{10\,\delta g,\,
 5\times10^{-10}\max(\omega/k_0,1)\right\}.
\]
The splitting and radial-shift log slopes must each lie within \(0.1\) of
\(1\).  The frequency-remainder slope must lie within \(0.2\) of \(2\).
Reversing \(\epsilon\) must exchange the axis and diagonal roles to angular
tolerance \(10^{-12}\).

\item The critical stage must return eight resolved alternating roots, four
minima and four saddles, whose Poincar\'e indices sum to zero.  This workflow
requirement does not claim exhaustive continuum enumeration.  Candidate-grid
doubling must preserve count and kind, while position, frequency, and Hessian
differences obey the uncertainty bounds stated above.  The negative-\(\epsilon\)
run must preserve angles within \(2\times10^{-6}\) and exchange
minimum/saddle roles.  Hessian separation, ten-times boundary separation,
coarse/fine winding agreement, the fine \(<\pi/2\) direction increment, and
annular index closure are all mandatory.

\item A full-wave amplitude must satisfy the order-pair inequality stated in
the response subsection.  At each anisotropy, the minimum relative branch gap
over all response surfaces must exceed ten times their maximum relative
frequency uncertainty.  A critical-point frequency
and its independently sampled modal-node frequency must differ by no more
than the node frequency uncertainty plus
\(64\epsilon_{\rm mach}\max(\omega_{\rm point},\omega_{\rm node},1)\).
The three-grid complex response, rate, and phase gates in the preceding
subsection are mandatory.

\item The Green workflow requires at least ten uniform-window samples.  Its
maximum Bessel-collapse error without fitted alignment is \(0.08\).  The early and late
slopes must be within \(0.05\) of \(-1/2\) and \(-1\), respectively, and the
last two grids' early and late slopes must differ by at most \(0.05\).  The
fixed-Morse relative RMS error, relative maximum error, and cancellation-region
normalized RMS must not exceed \(0.055\), \(0.17\), and \(0.09\).  The
crossover-time log slope must be within \(0.10\) of \(-1\), and both
production phase-discrepancy estimators must pass the configured \(0.05\)
acceptance threshold.

\item The auxiliary response study requires strictly decreasing quadrature
and interpolation errors, successive error ratios at most \(0.5\), finest
relative quadrature error at most \(0.05\), and the finest accumulated
interpolation phase-discrepancy estimator at most \(0.05\).  Its baseline source/window response
ratio must equal one with relative and absolute tolerances
\(5\times10^{-13}\), the robustness sweep must have nonzero range, and the
widest-window boundary weight must not exceed \(5\times10^{-3}\).

\item The silicon stress test must return eight points and pass the same
noncritical-boundary and annular-index closure checks.  It remains a
finite-anisotropy stress test, not a weak-anisotropy proof.
\end{enumerate}
}

At least ten samples are required in any direct log--log power-law fit;
late-time beat fits additionally require at least three samples per half-open
bin and at least four complete bins.  Changing any fixed threshold changes
the method and requires regenerating the artifact sidecars.

The committed full-profile discretizations are also fixed in the workflow
modules.  The isotropic stage uses GLL orders \(12/16\), 81 radial branch
nodes, and 41 local-curvature nodes.  The sensitivity stage uses order \(10\),
64 endpoint-free angles, constitutive differencing step \(0.25\), and the
independent physical finite-difference step \(0.0025\).  The weak-cubic
critical-point and scaling stages use nested orders \(10/14\); the
critical-point artifact uses \(\epsilon=0.02\), a \(26\times26\) display
grid, and Hessian step \(10^{-3}\).  The Green stage uses nested orders
\(10/14\), \(\epsilon\in\{0.005,0.01,0.02,0.04,0.08\}\), and
\(t_n=400+25n\) for \(0\leq n\leq504\).  The convergence stage sweeps
coarse/fine order pairs \(6/10,8/12,10/14,12/16,14/18\), uses an independent
\(513\times128\) response reference surface, and uses sensitivity steps
\(0.02,0.01,0.005,0.0025\).  The silicon stress test uses nested orders
\(12/16\), a \(40\times40\) display grid, and Hessian step
\(3\times10^{-3}\).  These are full-profile settings; the smoke profile is a
runtime check and does not replace any committed full-profile artifact.

The executable interface chain is explicit.  The first seven entries below
construct and track the spectral problem, the next three compute sensitivities
and critical points, and the final three certify and fit the response.  The
six trailing constants fix the Green-workflow masks.
{\footnotesize
\begin{verbatim}
gll_nodes_weights
differentiation_matrix
build_gll_mesh
assemble_plate_matrices
assemble_wavevector_derivatives
solve_plate_modes
mass_mac
RingAnchoredSpectralEvaluator
radial_frequency_sensitivity
locate_critical_points
verify_annular_exhaustion
integrate_registered_grid_convergence
assert_phase_accuracy
fit_power_law
_EARLY_TAU_WINDOW
_FIXED_EPSILON
_MORSE_COMPARISON_START
_MORSE_COMPARISON_STOP
_MORSE_COHERENCE_THRESHOLD
_UNIFORM_TAU_MAXIMUM
\end{verbatim}
}

\subsection{Qualified numerical-artifact registry}

The seven versioned artifacts are the sole numerical evidence consumed by
the figures and manuscript.  Each NPZ array is named as
\texttt{artifact.key}; each value in the JSON sidecar's historically named
\texttt{tolerances} namespace is
\texttt{artifact.tolerances.key}.  That namespace contains a mixture of
configured limits and measured diagnostics; a measured value is not itself a
pass threshold.  The hard inequalities are the list above.  The literal
registry below is complete for the committed schema and sidecars.  Array
meanings are grouped before it to make the mapping operational rather than
decorative.

\paragraph{Isotropic baseline.}
The isotropic artifact stores the sampled symmetric branch, exact ZGV
coordinates and curvature, its local quadratic comparison, and the selected
thickness mode and squared-displacement proxy.  Its sidecar binds exact/spectral mismatch,
nested-order frequency uncertainty, and eigengap.

\paragraph{Angular perturbation.}
The sensitivity artifact stores \(V(\theta)\), \(B(\theta)\), their
finite-difference checks, cubic harmonic coefficients, and the physical
signed shift \(\epsilon V_4\).  Its sidecar binds eigenpair, symmetry,
harmonic-leakage, and relative derivative errors.

\paragraph{Critical points and scaling.}
The critical-point artifact stores the eight full-wave roots, Cartesian
Hessian eigenvalues and Poincar\'e indices, perturbative predictions, and the
two plotted Cartesian surfaces.  The scaling artifact stores signed
\(\epsilon\) frequency and radial shifts, their compensated remainders, role
reversal, and measured log slopes.

\paragraph{Response, convergence, and stress test.}
The response artifact stores three complex responses: direct,
uniform-normal-form, and exact-Morse.  It also stores masks, spectra, slopes,
critical frequencies, and phase errors.  The convergence artifact holds the
polynomial, tracking, response-grid, finite-difference, and source/window
sweeps.  The silicon artifact is a finite-anisotropy stress test only.  It is
not evidence for the weak-anisotropy theorem.

{\footnotesize
\begin{verbatim}
isotropic_zgv.kappa
isotropic_zgv.omega_symmetric
isotropic_zgv.branch_labels
isotropic_zgv.kappa0
isotropic_zgv.omega0
isotropic_zgv.curvature_a
isotropic_zgv.local_q
isotropic_zgv.local_omega
isotropic_zgv.local_quadratic
isotropic_zgv.mode_z
isotropic_zgv.mode_u
isotropic_zgv.mode_squared_displacement
isotropic_zgv.tolerances.configured_isotropic_match
isotropic_zgv.tolerances.maximum_frequency_uncertainty
isotropic_zgv.tolerances.minimum_relative_eigengap
isotropic_zgv.tolerances.relative_isotropic_frequency

angular_sensitivity.theta
angular_sensitivity.V
angular_sensitivity.B
angular_sensitivity.V_reconstruction
angular_sensitivity.harmonic_order
angular_sensitivity.harmonic_amplitude
angular_sensitivity.V0
angular_sensitivity.V4
angular_sensitivity.V8
angular_sensitivity.epsilon
angular_sensitivity.delta_c
angular_sensitivity.physical_V4_shift
angular_sensitivity.V_fd
angular_sensitivity.B_fd
angular_sensitivity.tolerances.configured_relative_sensitivity
angular_sensitivity.tolerances.maximum_eigen_residual
angular_sensitivity.tolerances.minimum_relative_eigengap
angular_sensitivity.tolerances.mirror_defect
angular_sensitivity.tolerances.non_cubic_leakage
angular_sensitivity.tolerances.periodicity_defect
angular_sensitivity.tolerances.relative_B_error
angular_sensitivity.tolerances.relative_V_error

critical_points.kx
critical_points.ky
critical_points.kappa
critical_points.theta
critical_points.omega
critical_points.hessian_eigenvalues
critical_points.morse_index
critical_points.kind
critical_points.kx_pred
critical_points.ky_pred
critical_points.omega_pred
critical_points.gradient_residual
critical_points.kx_grid
critical_points.ky_grid
critical_points.omega_iso_grid
critical_points.omega_aniso_grid
critical_points.tolerances.epsilon
critical_points.tolerances.maximum_gradient_residual
critical_points.tolerances.maximum_grid_frequency_uncertainty
critical_points.tolerances.minimum_hessian_to_uncertainty
critical_points.tolerances.negative_candidate_grid_frequency_difference
critical_points.tolerances.negative_candidate_grid_hessian_difference
critical_points.tolerances.negative_candidate_grid_position_difference
critical_points.tolerances.negative_maximum_boundary_gradient_uncertainty
critical_points.tolerances.negative_minimum_boundary_gradient
critical_points.tolerances.positive_candidate_grid_frequency_difference
critical_points.tolerances.positive_candidate_grid_hessian_difference
critical_points.tolerances.positive_candidate_grid_position_difference
critical_points.tolerances.positive_maximum_boundary_gradient_uncertainty
critical_points.tolerances.positive_minimum_boundary_gradient

perturbation_scaling.epsilon
perturbation_scaling.delta_omega_full
perturbation_scaling.delta_omega_pred
perturbation_scaling.q_min_full
perturbation_scaling.q_min_pred
perturbation_scaling.q_saddle_full
perturbation_scaling.q_saddle_pred
perturbation_scaling.omega_min_error
perturbation_scaling.omega_saddle_error
perturbation_scaling.compensated_splitting
perturbation_scaling.compensated_q_min
perturbation_scaling.compensated_q_saddle
perturbation_scaling.compensated_frequency_error
perturbation_scaling.role_reversal_theta_min
perturbation_scaling.role_reversal_theta_saddle
perturbation_scaling.slope_splitting
perturbation_scaling.slope_radial
perturbation_scaling.slope_remainder
perturbation_scaling.tolerances.isotropic_axis_radius_bias
perturbation_scaling.tolerances.isotropic_diagonal_radius_bias
perturbation_scaling.tolerances.maximum_frequency_uncertainty
perturbation_scaling.tolerances.maximum_gradient_residual
perturbation_scaling.tolerances.slope_radial_error
perturbation_scaling.tolerances.slope_remainder_error
perturbation_scaling.tolerances.slope_splitting_error

green_crossover.epsilon
green_crossover.time
green_crossover.G_full
green_crossover.G_normal_form
green_crossover.G_morse
green_crossover.tau
green_crossover.J0
green_crossover.scaled_response
green_crossover.envelope
green_crossover.rms_envelope
green_crossover.local_slope
green_crossover.fit_window_early
green_crossover.fit_window_late
green_crossover.slope_early
green_crossover.slope_late
green_crossover.crossover_time
green_crossover.spectrum_omega
green_crossover.spectrum
green_crossover.omega_min
green_crossover.omega_saddle
green_crossover.phase_error
green_crossover.tolerances.bessel_collapse_error_by_epsilon
green_crossover.tolerances.configured_phase_error
green_crossover.tolerances.fixed_epsilon_cancellation_fraction
green_crossover.tolerances.fixed_epsilon_cancellation_region_normalized_rms
green_crossover.tolerances.fixed_epsilon_morse_relative_max_error
green_crossover.tolerances.fixed_epsilon_morse_relative_rms_error
green_crossover.tolerances.maximum_bessel_collapse_error
green_crossover.tolerances.maximum_phase_error
green_crossover.tolerances.maximum_relative_frequency_uncertainty_by_epsilon
green_crossover.tolerances.measured_crossover_log_slope
green_crossover.tolerances.measured_early_slope_error
green_crossover.tolerances.measured_late_slope_error
green_crossover.tolerances.minimum_anisotropic_branch_eigengap_by_epsilon
green_crossover.tolerances.registered_accumulated_phase_errors
green_crossover.tolerances.registered_complex_response_errors
green_crossover.tolerances.registered_early_slope_by_grid
green_crossover.tolerances.registered_late_slope_by_grid
green_crossover.tolerances.registered_nested_frequency_errors

convergence.polynomial_order
convergence.omega0_error
convergence.kappa0_error
convergence.curvature_error
convergence.eigen_residual
convergence.hermitian_residual
convergence.mass_orthogonality
convergence.eigengap
convergence.angular_resolution
convergence.radial_resolution
convergence.quadrature_error
convergence.interpolation_error
convergence.phase_error
convergence.sensitivity_step
convergence.V4_fd_error
convergence.B_fd_error
convergence.tracking_kappa
convergence.tracking_mac
convergence.tracking_gap
convergence.source_width
convergence.window_width
convergence.response_sensitivity
convergence.tolerances.configured_curvature_match
convergence.tolerances.configured_eigen_residual
convergence.tolerances.configured_isotropic_match
convergence.tolerances.configured_phase_error
convergence.tolerances.final_B_fd_relative_error
convergence.tolerances.final_V4_fd_relative_error
convergence.tolerances.final_relative_curvature_error
convergence.tolerances.final_relative_kappa0_error
convergence.tolerances.final_relative_omega0_error
convergence.tolerances.finest_accumulated_phase_error
convergence.tolerances.finest_interpolation_frequency_error
convergence.tolerances.finest_quadrature_relative_error
convergence.tolerances.maximum_eigen_residual
convergence.tolerances.maximum_hermitian_residual
convergence.tolerances.maximum_mass_orthogonality
convergence.tolerances.minimum_eigengap
convergence.tolerances.minimum_tracking_gap
convergence.tolerances.minimum_tracking_mac

silicon_stress_test.kx_grid
silicon_stress_test.ky_grid
silicon_stress_test.omega_grid
silicon_stress_test.kx
silicon_stress_test.ky
silicon_stress_test.omega
silicon_stress_test.hessian_eigenvalues
silicon_stress_test.kind
silicon_stress_test.tracking_gap
silicon_stress_test.gradient_residual
silicon_stress_test.tolerances.maximum_boundary_gradient_uncertainty
silicon_stress_test.tolerances.maximum_gradient_residual
silicon_stress_test.tolerances.minimum_boundary_gradient
silicon_stress_test.tolerances.minimum_tracking_gap
\end{verbatim}
}

\subsection{Reconstruction boundary}

An independent reader can reconstruct the calculation from the weak-form
matrices, the qualified configuration, and the listed interfaces.  The
numerical evidence is limited to the tracked local branch inside the
registered annulus and spectral window.  It does not claim completeness of
the full plate spectrum outside that window, and the finite-silicon run does
not extend the weak-anisotropy proof beyond its perturbative assumptions.
